% Gliederung Stand 30.10.2026. ~10 Seiten.
% Schema je Unterkapitel: Anforderung -> Alternativen -> begründete Entscheidung -> Kern der Lösung -> Verweis auf Nachweis in Kap. 8.
\chapter{Konzept der Umgebung}
\label{chap:konzept}
% Kurze Kapiteleinleitung vor 6.1.

\section{Ausgangsbasis und Abgrenzung}
\label{sec:ausgangsbasis}

\subsection{Auswahl der Ausgangsbasis}
\label{sec:auswahl_basis}
% g1pilot (HuCeBot) gegenüber Alternativen (Unitree-Beispiele, Neuentwicklung).

\subsection{Übernommener Funktionsumfang und Defizite}
\label{sec:defizite}
% Fehlte: Lauf-Simulation, Real-Sicherheit, Hände, Laienbedienung.

\subsection{Abgrenzung der Eigenleistung}
\label{sec:eigenleistung}
% Tabelle "übernommen aus g1pilot vs. neu entwickelt".

\section{Gesamtarchitektur}
\label{sec:architektur}
% [A2, F1]

\subsection{Komponenten und Datenflüsse}
\label{sec:komponenten}
% Kernbild des Kapitels.

\subsection{Umschaltung zwischen Simulation und Realbetrieb}
\label{sec:sim_real_schalter}
% Eine Umgebungsvariable, getrennte DDS-Domains.

\subsection{Fortbewegungspfade in Simulation und Realbetrieb}
\label{sec:fortbewegungspfade}
% Sim-Ersatzregler vs. Onboard-Regler hinter gleichen Topics.

\section{Sicherheitskonzept}
\label{sec:sicherheitskonzept}
% [A5, F3]

\subsection{Betriebszustände und Übergänge}
\label{sec:zustaende}
% Zustandsmaschine als Diagramm.

\subsection{Schutzmechanismen}
\label{sec:schutzmechanismen}
% Limits, Totmann, Notaus-Kette, Bodenkontaktprüfung.

\subsection{Übergabe der Arme an den Onboard-Regler}
\label{sec:armuebergabe}

\section{Bedienkonzept}
\label{sec:bedienkonzept}
% [A6, F3]

\subsection{Leitgedanken}
\label{sec:leitgedanken}
% Ein Modus zur Zeit (GEHEN/GREIFEN), Simulation zuerst.

\subsection{Bedienwege und geteilter Bedienzustand}
\label{sec:bedienwege_konzept}

\subsection{Posen und Abläufe als aufgabenunabhängige Bausteine}
\label{sec:posen_ablaeufe}
% [A1]

\section{Erweiterungskonzept}
\label{sec:erweiterungskonzept}
% [A7, F4]

\subsection{Andockpunkte für externe Projekte}
\label{sec:andockpunkte}
% Arm-API, DDS-Schnittstelle, Szenen, ROS-Topics.

\subsection{Modularisierung über Container-Profile}
\label{sec:container_profile}
